\documentclass[11pt]{article}

\usepackage[utf8]{inputenc}
\usepackage[T1]{fontenc}
\usepackage{lmodern}
\usepackage{geometry}
\usepackage{amsmath,amssymb,amsfonts,amsthm,mathtools}
\usepackage{bm}
\usepackage{enumitem}
\usepackage{microtype}
\usepackage{hyperref}
\usepackage{titlesec}
\usepackage{booktabs}
\usepackage{array}
\usepackage{setspace}
\usepackage{mathrsfs}
\usepackage{physics}

\geometry{margin=1in}
\hypersetup{colorlinks=true, linkcolor=black, urlcolor=blue, citecolor=black}
\setlist[enumerate]{topsep=0.4em,itemsep=0.35em}
\setlist[itemize]{topsep=0.3em,itemsep=0.25em}
\titleformat{\section}{\large\bfseries}{\thesection.}{0.5em}{}
\titleformat{\subsection}{\normalsize\bfseries}{\thesubsection.}{0.5em}{}
\setstretch{1.06}

\newcommand{\Aether}{\AE{}ther}
\newcommand{\Aflow}{\AE{}ther-flow}
\newcommand{\TheoryName}{\Aflow{} Interpretation of Relativity}
\newcommand{\TheoryProgram}{\Aether{} / \Aflow{} framework}

\newcommand{\CompletedMetric}{g^{\sharp}}

\newcommand{\Car}{\mathrm{car}}
\newcommand{\Cur}{\mathrm{cur}}
\newcommand{\Hom}{\mathrm{hom}}

\newcommand{\ForSpace}[1]{\mathcal I_{#1}^{\mathrm{for}}}
\newcommand{\PrimShell}{\mathcal S_{\mathrm{prim}}}

\newcommand{\Reduction}[1]{\mathcal R_{#1}}
\newcommand{\CurrentCarrier}[1]{\mathfrak C_{#1}^{\mathrm{cur}}}
\newcommand{\EqCarrier}[1]{\mathfrak C_{#1}^{\mathrm{eq}}}
\newcommand{\CurrentExtract}[1]{\mathscr E_{#1}^{\mathrm{cur}}}
\newcommand{\EqExtract}[1]{\mathscr E_{#1}^{\mathrm{eq}}}
\newcommand{\PrimCurrent}[1]{\mathcal J_{#1}^{\mathrm{prim}}}
\newcommand{\PrimTheta}[1]{\Theta_{#1}^{\mathrm{prim}}}

\newcommand{\MetricOut}{\mathscr G_{\mathrm{out}}}
\newcommand{\MatterOut}{\mathscr T_{\mathrm{out}}}

\newcommand{\CoordSpace}[1]{Q_{#1}}
\newcommand{\CoordBody}[1]{K_{#1}}
\newcommand{\NormalSpace}[1]{N_{#1}}
\newcommand{\NormalBody}[1]{B_{#1}}
\newcommand{\VisibleAffine}[1]{\widehat X_{#1}}
\newcommand{\AffineData}[1]{\widehat{\mathcal D}_{#1}}
\newcommand{\ConstraintTarget}[1]{E_{#1}}
\newcommand{\ConstraintMap}[1]{C_{#1}}
\newcommand{\ConstraintValue}[1]{c_{#1}}
\newcommand{\NormalOperator}[1]{H_{#1}}
\newcommand{\SymGroup}[1]{G_{#1}}
\newcommand{\TimeRev}[1]{\tau_{#1}}
\newcommand{\ShellChart}[1]{\chi_{#1}}

\newcommand{\ShellAffine}[1]{A_{#1}^{\mathrm{sh}}}
\newcommand{\ShellDir}[1]{V_{#1}^{\mathrm{sh}}}
\newcommand{\ShellNormal}[1]{N_{#1}^{\mathrm{sh}}}
\newcommand{\ShellProj}[1]{P_{#1}^{\mathrm{sh},\perp}}
\newcommand{\ShellOffset}[1]{b_{#1}^{\mathrm{sh}}}
\newcommand{\ShellEmbed}[1]{\iota_{#1}^{\mathrm{sh}}}
\newcommand{\AffineNucleus}[1]{\mathfrak N_{#1}^{\mathrm{aff}}}

\newtheorem{definition}{Definition}
\newtheorem{lemma}{Lemma}
\newtheorem{proposition}{Proposition}
\newtheorem{remark}{Remark}
\newtheorem{corollary}{Corollary}
\newtheorem{theorem}{Theorem}

\title{\textbf{The \TheoryName{}}\\[0.4em]
\large Primitive-Reservoir Observer-Reduction\\
\large Primitive Reversible Constrained Dynamics\\
\large Collapse to a Minimal Exact-Shell Affine Nucleus}
\author{}
\date{}

\begin{document}

\maketitle

\begin{abstract}
The current deepest benchmark-facing note on the active primitive-reservoir observer-reduction line already made the decisive bridge-object forcing move now on record: the primitive reversible constrained substrate dynamics normal form is no longer merely available, because it is derived canonically from an ambient shell-law substrate package with hidden-sector gap and is necessary there up to the forced translated-shell and orthogonal-hidden coordinate identifications. The remaining honest bridge-object question is therefore no longer whether that larger normal form exists. The remaining question is whether every constituent still listed in that full normal form is independently benchmark-facing, or whether a sharper theorem-level collapse is available strictly inside it.

The present note takes that sharper internal route. For each sector it isolates a smaller \emph{exact-shell affine nucleus}. The nucleus keeps the shell-coordinate body, the affine shell-data readout, the exact shell affine plane together with the exact shell realization of the primitive forcing shell, the off-shell body and positive off-shell operator, the induced symmetry and time-reversal structure, the shell solvability property, and the affine-kernel coercivity package. What it omits as independent primitive data are the explicit constraint-target space, the linear constraint map, and the fixed constraint value. Those are shown to be canonically induced by orthogonal projection onto the normal complement of the shell affine plane.

This is a genuine bridge-object sharpening rather than another deeper same-output relay. Once the exact-shell affine nucleus is fixed, the larger primitive reversible constrained dynamics normal form is recovered canonically, and any compatible full presentation differs only by an inessential Euclidean presentation of the same shell-normal target. Consequently all current benchmark-facing downstream uses of the full normal form, including the explicit substrate equation / symmetry / constraint package, the unique selector, and the same benchmark one-metric observer package, already factor through the smaller nucleus. The benchmark boundary remains fixed throughout: the operative metric is still exactly \(\CompletedMetric\), the ordinary matter route is unchanged, there is no second operative metric, and the low-energy matter law is not enlarged. The open burden therefore sharpens once more. The next honest primary manuscript must derive or justify this smaller exact-shell affine nucleus from explicit substrate structure in a benchmark-facing way, or else prove a still smaller theorem-level collapse strictly inside it.
\end{abstract}

\section{Introduction}

The active derivational program of \TheoryName{} has now reached the point where the bridge object itself must be sharpened rather than merely moved one layer deeper again. The observer-reduction datum line, the defect-fiber factorization theorem, the one-metric output theorem, the chain-forcing theorem, the selector-existence theorem, the stationary-construction theorem, the explicit-package necessity / uniqueness theorem, and the later primitive-normal-form forcing theorem have already discharged what was honestly open at their respective levels. The present routing layer matters because it changes what now counts as real progress. Once the primitive reversible constrained substrate dynamics normal form is itself forced from an ambient shell law, another manuscript that merely inserts a deeper same-output relay below that larger normal form no longer addresses the live burden.

There are two disciplined ways to continue. One can derive or justify the bridge object itself from still more primitive explicit substrate structure in a way that changes benchmark-facing status. Or one can prove that the current bridge object still contains redundant presentation freedom and collapse it to a smaller theorem-level datum strictly inside the same object. The present note takes that second route.

Its gain is narrower than a completed first-principles substrate derivation and stronger than another same-output deeper relay. The key observation is simple. In the current primitive reversible constrained dynamics normal form, the exact shell is already presented as an affine slice of the shell-coordinate body. Once that affine shell plane, the shell-data readout, the exact shell realization, the off-shell sector, the reversible structure, and the solvability / coercive package are fixed, the separate linear constraint target, linear constraint map, and fixed constraint value are no longer extra benchmark-facing freedom. They are canonically induced by the Euclidean orthogonal decomposition determined by that affine shell plane.

\paragraph{Framework and claim boundary.}
\TheoryProgram{} denotes the broader \Aether{} / \Aflow{} framework built on the claims that the \Aether{} is the underlying four-dimensional substrate of reality, the \Aflow{} is its intrinsic ordered motion, observed three-dimensional space is the local experiential slice of that deeper substrate, S-time is the experienced order of change arising from matter, light, and the \Aflow{}, and observed expansion is the three-dimensional appearance of deeper four-dimensional ordered motion. Within that framework, \TheoryName{} denotes the exact-closure theory in which the effective gravitational dynamics are adopted to be exactly Einsteinian with universal matter coupling. In the active sequence, \emph{adoption} means use of that established relativistic dynamics without claiming substrate derivation, while \emph{derivation} is reserved for a first-principles recovery from explicit substrate variables.

The present note stays strictly inside that boundary. It does not alter the benchmark exact-closure package. It does not introduce a second operative metric or a new low-energy matter law. It only proves that the larger primitive reversible constrained dynamics normal form collapses further to a smaller exact-shell affine nucleus at current bridge-object scope.

\section{Fixed Primitive Continuation Data}

\begin{definition}[Recorded primitive continuation data]
\label{def:fixed}
Fix the following continuation data from the active primitive-reservoir observer-reduction line.
\begin{enumerate}
    \item For each sector label \(\sigma\in\{\Car,\Cur,\Hom\}\), a primitive forcing shell
    \[
    \ForSpace{\sigma},
    \]
    a visible reduction map
    \[
    \Reduction{\sigma}:\ForSpace{\sigma}\to X_{\sigma},
    \]
    and shell-level current and equilibrium carriers
    \[
    \CurrentCarrier{\sigma}:\ForSpace{\sigma}\to Z_{\sigma}^{\mathrm{cur}},
    \qquad
    \EqCarrier{\sigma}:\ForSpace{\sigma}\to Z_{\sigma}^{\mathrm{eq}}.
    \]
    \item The product shell
    \[
    \PrimShell
    \equiv
    \ForSpace{\Car}\times \ForSpace{\Cur}\times \ForSpace{\Hom}.
    \]
    \item Fixed shell extractors
    \[
    \CurrentExtract{\sigma}:Z_{\sigma}^{\mathrm{cur}}\to Y_{\sigma}^{\mathrm{cur}},
    \qquad
    \EqExtract{\sigma}:Z_{\sigma}^{\mathrm{eq}}\to Y_{\sigma}^{\mathrm{eq}},
    \]
    together with the recorded shell composites
    \begin{equation}
    \PrimCurrent{\sigma}
    =
    \CurrentExtract{\sigma}\circ \CurrentCarrier{\sigma},
    \qquad
    \PrimTheta{\sigma}
    =
    \EqExtract{\sigma}\circ \EqCarrier{\sigma}.
    \label{eq:primcomposites}
    \end{equation}
    \item The recorded metric and ordinary-matter outputs
    \[
    \MetricOut:\PrimShell\to\{\text{observer metrics}\},
    \qquad
    \MatterOut:\PrimShell\times\Psi\to\{\text{observer stress tensors}\},
    \]
    satisfying
    \begin{equation}
    \MetricOut(\mathbf w)=\CompletedMetric,
    \qquad
    \MatterOut(\mathbf w,\psi)
    =
    T_{\mu\nu}^{\mathrm{matter}}[\CompletedMetric,\psi]
    \label{eq:fixedoutputs}
    \end{equation}
    for every \(\mathbf w\in\PrimShell\) and every matter configuration \(\psi\in\Psi\).
\end{enumerate}
\end{definition}

\begin{remark}
Definition \ref{def:fixed} is not reopened here. The primitive shell data, the one operative metric, and the ordinary-matter route remain fixed continuation data. The present note addresses only the benchmark-facing internal compression of the current bridge object.
\end{remark}

\section{Recorded Full Primitive Reversible Constrained Dynamics}

\begin{definition}[Primitive reversible constrained dynamics normal form]
\label{def:full}
Fix \(\sigma\in\{\Car,\Cur,\Hom\}\). A \emph{primitive reversible constrained dynamics normal form} for sector \(\sigma\) consists of the following data.
\begin{enumerate}
    \item Finite-dimensional Euclidean spaces
    \[
    \CoordSpace{\sigma},
    \qquad
    \NormalSpace{\sigma},
    \qquad
    \VisibleAffine{\sigma},
    \qquad
    \ConstraintTarget{\sigma},
    \]
    together with compact convex bodies
    \[
    \CoordBody{\sigma}\subset\CoordSpace{\sigma},
    \qquad
    \NormalBody{\sigma}\subset\NormalSpace{\sigma},
    \]
    each with nonempty interior, and with \(0\in\operatorname{int}\NormalBody{\sigma}\).
    \item A smooth embedding
    \[
    \jmath_{\sigma}:X_{\sigma}\hookrightarrow \VisibleAffine{\sigma}.
    \]
    \item An affine shell-data readout
    \[
    \AffineData{\sigma}:\CoordSpace{\sigma}\to
    \VisibleAffine{\sigma}\times Z_{\sigma}^{\mathrm{cur}}\times Z_{\sigma}^{\mathrm{eq}}
    \]
    whose image on \(\CoordBody{\sigma}\) lies in
    \[
    \jmath_{\sigma}(X_{\sigma})\times Z_{\sigma}^{\mathrm{cur}}\times Z_{\sigma}^{\mathrm{eq}}.
    \]
    \item A linear constraint map
    \[
    \ConstraintMap{\sigma}:\CoordSpace{\sigma}\to \ConstraintTarget{\sigma}
    \]
    and a fixed constraint value
    \[
    \ConstraintValue{\sigma}\in\ConstraintTarget{\sigma}.
    \]
    \item A positive self-adjoint operator
    \[
    \NormalOperator{\sigma}:\NormalSpace{\sigma}\to\NormalSpace{\sigma}
    \]
    with lower bound
    \begin{equation}
    \ev{\eta,\NormalOperator{\sigma}\eta}\ge \mu_{\sigma}\,\|\eta\|^{2}
    \qquad
    \text{for all }\eta\in\NormalSpace{\sigma}
    \label{eq:fullgap}
    \end{equation}
    for some \(\mu_{\sigma}>0\).
    \item A compact symmetry group
    \[
    \SymGroup{\sigma}\subset O(\CoordSpace{\sigma})\times O(\NormalSpace{\sigma})
    \]
    and an orthogonal involution
    \[
    \TimeRev{\sigma}\in O(\CoordSpace{\sigma})\times O(\NormalSpace{\sigma})
    \]
    preserving \(\CoordBody{\sigma}\), \(\NormalBody{\sigma}\), \(\AffineData{\sigma}\), \(\ConstraintMap{\sigma}\), \(\ConstraintValue{\sigma}\), and \(\NormalOperator{\sigma}\).
    \item A shell chart
    \[
    \ShellChart{\sigma}:\ForSpace{\sigma}\to
    \CoordBody{\sigma}\cap \ConstraintMap{\sigma}^{-1}(\ConstraintValue{\sigma})
    \]
    that is a diffeomorphism onto the full constrained shell-coordinate locus and satisfies
    \begin{equation}
    \AffineData{\sigma}\bigl(\ShellChart{\sigma}(w)\bigr)
    =
    \bigl(
    \jmath_{\sigma}(\Reduction{\sigma}(w)),
    \CurrentCarrier{\sigma}(w),
    \EqCarrier{\sigma}(w)
    \bigr)
    \label{eq:fullcompat}
    \end{equation}
    for every \(w\in\ForSpace{\sigma}\).
    \item Data-preserving shell solvability:
    \begin{equation}
    \AffineData{\sigma}(\CoordBody{\sigma})
    =
    \AffineData{\sigma}
    \bigl(
    \CoordBody{\sigma}\cap\ConstraintMap{\sigma}^{-1}(\ConstraintValue{\sigma})
    \bigr).
    \label{eq:fullsolvability}
    \end{equation}
    \item Data-kernel coercivity: there exists \(\kappa_{\sigma}>0\) such that for every
    \[
    \delta q\in \ker D\AffineData{\sigma},
    \qquad
    \delta n\in \NormalSpace{\sigma},
    \]
    one has
    \begin{equation}
    \|\ConstraintMap{\sigma}\delta q\|^{2}
    +
    \ev{\delta n,\NormalOperator{\sigma}\delta n}
    \ge
    \kappa_{\sigma}\,
    \bigl(
    \|\delta q\|^{2}+\|\delta n\|^{2}
    \bigr).
    \label{eq:fullcoercive}
    \end{equation}
\end{enumerate}
\end{definition}

\begin{remark}
Definition \ref{def:full} is the current bridge-object package whose full benchmark-facing status is at issue. The previous forcing theorem showed that this larger normal form is no longer merely available on the active line. The present note asks the sharper internal question: which constituents of that larger package remain genuinely benchmark-facing once one retains the exact shell affine geometry itself.
\end{remark}

\section{Minimal Exact-Shell Affine Nucleus}

\begin{definition}[Exact-shell affine nucleus]
\label{def:nucleus}
Fix \(\sigma\in\{\Car,\Cur,\Hom\}\). An \emph{exact-shell affine nucleus} \(\AffineNucleus{\sigma}\) for sector \(\sigma\) consists of the following data.
\begin{enumerate}
    \item Finite-dimensional Euclidean spaces
    \[
    \CoordSpace{\sigma},
    \qquad
    \NormalSpace{\sigma},
    \qquad
    \VisibleAffine{\sigma},
    \]
    together with compact convex bodies
    \[
    \CoordBody{\sigma}\subset\CoordSpace{\sigma},
    \qquad
    \NormalBody{\sigma}\subset\NormalSpace{\sigma},
    \]
    each with nonempty interior, and with \(0\in\operatorname{int}\NormalBody{\sigma}\).
    \item A smooth embedding
    \[
    \jmath_{\sigma}:X_{\sigma}\hookrightarrow \VisibleAffine{\sigma}.
    \]
    \item An affine shell-data readout
    \[
    \AffineData{\sigma}:\CoordSpace{\sigma}\to
    \VisibleAffine{\sigma}\times Z_{\sigma}^{\mathrm{cur}}\times Z_{\sigma}^{\mathrm{eq}}
    \]
    whose image on \(\CoordBody{\sigma}\) lies in
    \[
    \jmath_{\sigma}(X_{\sigma})\times Z_{\sigma}^{\mathrm{cur}}\times Z_{\sigma}^{\mathrm{eq}}.
    \]
    \item An affine shell plane
    \[
    \ShellAffine{\sigma}\subset\CoordSpace{\sigma}
    \]
    with translation space
    \[
    \ShellDir{\sigma}\equiv T\ShellAffine{\sigma}
    \]
    and shell-normal complement
    \[
    \ShellNormal{\sigma}\equiv \bigl(\ShellDir{\sigma}\bigr)^{\perp}.
    \]
    Let
    \[
    \ShellProj{\sigma}:\CoordSpace{\sigma}\to\ShellNormal{\sigma}
    \]
    denote the orthogonal projection.
    \item An exact shell realization
    \[
    \ShellEmbed{\sigma}:\ForSpace{\sigma}\to
    \CoordBody{\sigma}\cap \ShellAffine{\sigma}
    \]
    that is a diffeomorphism onto the full affine shell slice and satisfies
    \begin{equation}
    \AffineData{\sigma}\bigl(\ShellEmbed{\sigma}(w)\bigr)
    =
    \bigl(
    \jmath_{\sigma}(\Reduction{\sigma}(w)),
    \CurrentCarrier{\sigma}(w),
    \EqCarrier{\sigma}(w)
    \bigr)
    \label{eq:nucleuscompat}
    \end{equation}
    for every \(w\in\ForSpace{\sigma}\).
    \item A positive self-adjoint operator
    \[
    \NormalOperator{\sigma}:\NormalSpace{\sigma}\to\NormalSpace{\sigma}
    \]
    with lower bound
    \begin{equation}
    \ev{\eta,\NormalOperator{\sigma}\eta}\ge \mu_{\sigma}\,\|\eta\|^{2}
    \qquad
    \text{for all }\eta\in\NormalSpace{\sigma}
    \label{eq:nucleusgap}
    \end{equation}
    for some \(\mu_{\sigma}>0\).
    \item A compact symmetry group
    \[
    \SymGroup{\sigma}\subset O(\CoordSpace{\sigma})\times O(\NormalSpace{\sigma})
    \]
    and an orthogonal involution
    \[
    \TimeRev{\sigma}\in O(\CoordSpace{\sigma})\times O(\NormalSpace{\sigma})
    \]
    preserving \(\CoordBody{\sigma}\), \(\NormalBody{\sigma}\), \(\AffineData{\sigma}\), \(\ShellAffine{\sigma}\), and \(\NormalOperator{\sigma}\).
    \item Data-preserving shell solvability:
    \begin{equation}
    \AffineData{\sigma}(\CoordBody{\sigma})
    =
    \AffineData{\sigma}
    \bigl(
    \CoordBody{\sigma}\cap\ShellAffine{\sigma}
    \bigr).
    \label{eq:nucleussolvability}
    \end{equation}
    \item Affine-kernel coercivity: there exists \(\kappa_{\sigma}>0\) such that for every
    \[
    \delta q\in \ker D\AffineData{\sigma},
    \qquad
    \delta n\in \NormalSpace{\sigma},
    \]
    one has
    \begin{equation}
    \|\ShellProj{\sigma}\delta q\|^{2}
    +
    \ev{\delta n,\NormalOperator{\sigma}\delta n}
    \ge
    \kappa_{\sigma}\,
    \bigl(
    \|\delta q\|^{2}+\|\delta n\|^{2}
    \bigr).
    \label{eq:nucleuscoercive}
    \end{equation}
\end{enumerate}
\end{definition}

\begin{remark}
Definition \ref{def:nucleus} omits exactly the explicit constraint-target space, linear constraint map, and fixed constraint value as independent primitive data. The point of the present theorem is that, once the shell affine plane itself is kept, that omitted triple is no longer independently benchmark-facing.
\end{remark}

\section{Canonical Recovery of the Constraint Presentation}

\begin{proposition}[The shell affine plane canonically determines the constraint presentation]
\label{prop:constraint}
Assume Definition \ref{def:nucleus}. Then there exists a unique vector
\[
\ShellOffset{\sigma}\in\ShellNormal{\sigma}
\]
such that
\begin{equation}
\ShellAffine{\sigma}
=
\left\{
q\in\CoordSpace{\sigma}:
\ShellProj{\sigma}(q)=\ShellOffset{\sigma}
\right\}.
\label{eq:affineplane}
\end{equation}
Consequently, the canonical assignments
\begin{equation}
\ConstraintTarget{\sigma}^{\mathrm{can}}
\equiv
\ShellNormal{\sigma},
\qquad
\ConstraintMap{\sigma}^{\mathrm{can}}
\equiv
\ShellProj{\sigma},
\qquad
\ConstraintValue{\sigma}^{\mathrm{can}}
\equiv
\ShellOffset{\sigma}
\label{eq:canonicalconstraint}
\end{equation}
recover the exact shell slice as
\begin{equation}
\CoordBody{\sigma}\cap\ShellAffine{\sigma}
=
\CoordBody{\sigma}\cap
\bigl(\ConstraintMap{\sigma}^{\mathrm{can}}\bigr)^{-1}
\bigl(\ConstraintValue{\sigma}^{\mathrm{can}}\bigr).
\label{eq:sameplane}
\end{equation}
Moreover \(\SymGroup{\sigma}\) and \(\TimeRev{\sigma}\) preserve
\[
\ConstraintMap{\sigma}^{\mathrm{can}}
\quad\text{and}\quad
\ConstraintValue{\sigma}^{\mathrm{can}}.
\]
\end{proposition}

\begin{proof}
Because \(\ShellAffine{\sigma}\) is an affine plane parallel to \(\ShellDir{\sigma}\), every two points \(q,q'\in\ShellAffine{\sigma}\) differ by an element of \(\ShellDir{\sigma}\). Hence their orthogonal projections to \(\ShellNormal{\sigma}\) agree, because \(\ShellProj{\sigma}(q-q')=0\). This common value is the unique vector \(\ShellOffset{\sigma}\) of \eqref{eq:affineplane}. The identity \eqref{eq:sameplane} is then immediate.

Since \(\SymGroup{\sigma}\) and \(\TimeRev{\sigma}\) preserve \(\ShellAffine{\sigma}\), they preserve its translation space \(\ShellDir{\sigma}\), its orthogonal complement \(\ShellNormal{\sigma}\), and the orthogonal projection \(\ShellProj{\sigma}\). They also preserve the affine plane itself, hence preserve the unique shell-offset vector \(\ShellOffset{\sigma}\). Therefore the canonical constraint presentation is symmetry- and time-reversal-compatible.
\end{proof}

\begin{proposition}[The exact-shell affine nucleus canonically induces the full normal form]
\label{prop:recover}
Assume Definition \ref{def:nucleus}. Then the data
\[
\bigl(
\CoordSpace{\sigma},
\NormalSpace{\sigma},
\VisibleAffine{\sigma},
\ConstraintTarget{\sigma}^{\mathrm{can}},
\CoordBody{\sigma},
\NormalBody{\sigma},
\jmath_{\sigma},
\AffineData{\sigma},
\ConstraintMap{\sigma}^{\mathrm{can}},
\ConstraintValue{\sigma}^{\mathrm{can}},
\NormalOperator{\sigma},
\SymGroup{\sigma},
\TimeRev{\sigma},
\ShellChart{\sigma}\equiv\ShellEmbed{\sigma}
\bigr)
\]
define a primitive reversible constrained dynamics normal form in the sense of Definition \ref{def:full}.
\end{proposition}

\begin{proof}
Items (1), (2), (3), (5), (6), and (7) of Definition \ref{def:full} are already included in Definition \ref{def:nucleus}, with \(\ShellChart{\sigma}\) taken to be \(\ShellEmbed{\sigma}\). Item (4) is Proposition \ref{prop:constraint}. Item (8) is exactly \eqref{eq:nucleussolvability} together with \eqref{eq:sameplane}. Item (9) is \eqref{eq:nucleuscoercive}, again using \(\ConstraintMap{\sigma}^{\mathrm{can}}=\ShellProj{\sigma}\). Thus every exact-shell affine nucleus canonically induces a full primitive reversible constrained dynamics normal form.
\end{proof}

\begin{proposition}[Any compatible full presentation differs only by shell-normal target presentation]
\label{prop:uniqueness}
Let a full primitive reversible constrained dynamics presentation in the sense of Definition \ref{def:full} be \emph{compatible} with the exact-shell affine nucleus of Definition \ref{def:nucleus} if it has the same shell-coordinate space and body, the same off-shell space and body, the same shell-data readout, the same shell affine plane and exact shell realization, the same positive off-shell operator, the same symmetry / time-reversal structure, and the same solvability / coercive package. Then every such compatible full presentation satisfies:
\begin{enumerate}
    \item its constrained shell locus is exactly \(\CoordBody{\sigma}\cap\ShellAffine{\sigma}\);
    \item its constraint map has kernel \(\ShellDir{\sigma}\);
    \item there exists an isometric embedding
    \[
    U_{\sigma}:\ShellNormal{\sigma}\hookrightarrow \ConstraintTarget{\sigma}
    \]
    with image \(\operatorname{im}\ConstraintMap{\sigma}\) such that
    \begin{equation}
    \ConstraintMap{\sigma}
    =
    U_{\sigma}\circ \ShellProj{\sigma},
    \qquad
    \ConstraintValue{\sigma}
    =
    U_{\sigma}\bigl(\ShellOffset{\sigma}\bigr).
    \label{eq:targetpresentation}
    \end{equation}
\end{enumerate}
Hence the only remaining freedom in the larger full normal form is the inessential Euclidean presentation of the same shell-normal target.
\end{proposition}

\begin{proof}
Item (1) is the definition of compatibility. For item (2), if \(v\in\ShellDir{\sigma}\), then \(q+v\in\ShellAffine{\sigma}\) for every \(q\in\ShellAffine{\sigma}\), so by item (1) both \(q\) and \(q+v\) lie in the same constrained shell affine slice. Hence \(\ConstraintMap{\sigma}(v)=0\), proving \(\ShellDir{\sigma}\subset\ker\ConstraintMap{\sigma}\). Conversely, if \(v\in\ker\ConstraintMap{\sigma}\) and \(q\in\ShellAffine{\sigma}\), then
\[
\ConstraintMap{\sigma}(q+v)
=
\ConstraintMap{\sigma}(q)
=
\ConstraintValue{\sigma},
\]
so \(q+v\in\ShellAffine{\sigma}\). Therefore \(v\in\ShellDir{\sigma}\), and \(\ker\ConstraintMap{\sigma}=\ShellDir{\sigma}\).

For item (3), both \(\ConstraintMap{\sigma}\) and \(\ShellProj{\sigma}\) vanish exactly on \(\ShellDir{\sigma}\), so each factors through the quotient \(\CoordSpace{\sigma}/\ShellDir{\sigma}\). The Euclidean orthogonal decomposition
\[
\CoordSpace{\sigma}=\ShellDir{\sigma}\oplus\ShellNormal{\sigma}
\]
identifies that quotient with \(\ShellNormal{\sigma}\), and \(\ShellProj{\sigma}\) is the canonical quotient map in these coordinates. The image of \(\ConstraintMap{\sigma}\) is therefore an isometric copy of \(\ShellNormal{\sigma}\) inside \(\ConstraintTarget{\sigma}\), yielding an embedding \(U_{\sigma}\) such that \eqref{eq:targetpresentation} holds. The value statement follows because every point of \(\ShellAffine{\sigma}\) is sent to \(\ConstraintValue{\sigma}\) by \(\ConstraintMap{\sigma}\) and to \(\ShellOffset{\sigma}\) by \(\ShellProj{\sigma}\).
\end{proof}

\section{Downstream Bridge-Object Uses Factor Through the Nucleus}

\begin{proposition}[All current bridge-object uses factor through the exact-shell affine nucleus]
\label{prop:factor}
Assume Definition \ref{def:nucleus}. Then all current benchmark-facing downstream uses of the full primitive reversible constrained dynamics normal form already factor through the exact-shell affine nucleus together with the fixed primitive continuation data. In particular:
\begin{enumerate}
    \item the explicit substrate equation / symmetry / constraint package determined by the full normal form is canonically fixed by the nucleus;
    \item the sectorwise selector and the total selector determined at full-normal-form scope are the same for every compatible full presentation;
    \item the one operative metric \(\CompletedMetric\) and the same ordinary-matter route remain fixed.
\end{enumerate}
No separate larger constraint-target presentation is required for those benchmark-facing downstream uses once the nucleus is fixed.
\end{proposition}

\begin{proof}
By Proposition \ref{prop:recover}, the nucleus canonically induces a full primitive reversible constrained dynamics normal form. By Proposition \ref{prop:uniqueness}, every compatible full presentation differs only by the inessential Euclidean presentation of the same shell-normal target. The already recorded downstream constructions of the active line at full-normal-form scope use only the induced shell slice, shell-data readout, off-shell gap package, reversible structure, solvability, and coercive package. Those data are exactly the constituents of the nucleus. Therefore the explicit substrate package, the selector construction, and the benchmark one-metric outputs depend only on the nucleus together with the fixed primitive continuation data.
\end{proof}

\begin{proposition}[The benchmark package remains fixed]
\label{prop:benchmark}
Every exact-shell affine nucleus preserves the same benchmark one-metric package \eqref{eq:fixedoutputs}. In particular, the operative metric remains exactly \(\CompletedMetric\), the ordinary matter route is unchanged, there is no second operative metric, and the low-energy matter law is not enlarged.
\end{proposition}

\begin{proof}
Proposition \ref{prop:factor} shows that the benchmark-facing downstream outputs of the full primitive reversible constrained dynamics normal form already factor through the nucleus. Those outputs are fixed by Definition \ref{def:fixed}. Hence the benchmark package remains exactly the same.
\end{proof}

\section{Minimality of the Nucleus}

\begin{proposition}[Every nucleus constituent is benchmark-facing]
\label{prop:minimal}
No proper subtuple of the exact-shell affine nucleus \(\AffineNucleus{\sigma}\) still determines the full primitive reversible constrained dynamics normal form up to the compatible presentation freedom of Proposition \ref{prop:uniqueness}. More precisely:
\begin{enumerate}
    \item if \(\CoordBody{\sigma}\) is omitted, the shell-coordinate domain and hence the full shell slice are no longer fixed;
    \item if \(\AffineData{\sigma}\) is omitted, the visible / current / equilibrium shell readout is no longer fixed;
    \item if \(\ShellAffine{\sigma}\) is omitted, the induced constraint-target space, linear constraint map, and fixed constraint value are no longer fixed;
    \item if \(\ShellEmbed{\sigma}\) is omitted, the exact realization of the primitive forcing shell inside the shell slice is no longer fixed;
    \item if \(\NormalBody{\sigma}\) or \(\NormalOperator{\sigma}\) is omitted, the off-shell gap package is no longer fixed;
    \item if \(\SymGroup{\sigma}\) or \(\TimeRev{\sigma}\) is omitted, the reversible symmetry discipline of the bridge object is no longer fixed;
    \item if \eqref{eq:nucleussolvability} or \eqref{eq:nucleuscoercive} is omitted, the selector-supporting solvability / uniqueness package is no longer fixed.
\end{enumerate}
\end{proposition}

\begin{proof}
Each item names a distinct part of the reconstruction of Proposition \ref{prop:recover} and the downstream factorization of Proposition \ref{prop:factor}. Item (1) fixes where the shell affine plane is actually cut by admissible shell coordinates. Item (2) fixes the observer-relevant visible / current / equilibrium content. Item (3) is exactly the datum from which the omitted constraint presentation is reconstructed. Item (4) fixes the exact primitive forcing-shell realization required by \eqref{eq:nucleuscompat}. Item (5) fixes the off-shell gap data. Item (6) fixes the reversible symmetry package. Item (7) fixes the solvability and coercive mechanism that support the already-recorded selector and uniqueness consequences. Removing any one of these data destroys a distinct benchmark-facing part of the recovered full normal form.
\end{proof}

\section{Main Theorem}

\begin{theorem}[Primitive reversible constrained dynamics collapse to a minimal exact-shell affine nucleus]
\label{thm:main}
Fix the primitive continuation data of Definition \ref{def:fixed}. Then, for each sector \(\sigma\in\{\Car,\Cur,\Hom\}\), the larger primitive reversible constrained dynamics normal form of Definition \ref{def:full} collapses strictly inside itself to the smaller exact-shell affine nucleus \(\AffineNucleus{\sigma}\) of Definition \ref{def:nucleus}. More precisely:
\begin{enumerate}
    \item every full presentation in the sense of Definition \ref{def:full} canonically determines an exact-shell affine nucleus by retaining the shell-coordinate body, shell-data readout, exact shell affine plane, exact shell realization, off-shell gap package, reversible structure, and solvability / coercive package, and by forgetting only the explicit constraint-target space, linear constraint map, and fixed constraint value as independent data;
    \item the full primitive reversible constrained dynamics normal form is completely determined by the exact-shell affine nucleus, because the omitted constraint presentation is canonically recovered by Proposition \ref{prop:constraint};
    \item any compatible full presentation differs from the canonical recovered one only by the inessential Euclidean presentation of the shell-normal target described in Proposition \ref{prop:uniqueness};
    \item all current benchmark-facing downstream uses of the full primitive reversible constrained dynamics normal form, including the explicit substrate package, the unique selector, and the same benchmark one-metric observer package, already factor through the exact-shell affine nucleus by Proposition \ref{prop:factor};
    \item no proper subtuple of the exact-shell affine nucleus still determines the full primitive reversible constrained dynamics normal form up to the compatible presentation freedom of item (3).
\end{enumerate}
Hence the live bridge-object burden is no longer the larger full primitive reversible constrained dynamics normal form as such. What remains at current scope is the smaller minimal exact-shell affine nucleus.
\end{theorem}

\begin{proof}
Item (1) is the construction of Definition \ref{def:nucleus} read inside the larger package of Definition \ref{def:full}. Item (2) is Proposition \ref{prop:recover}, built on Proposition \ref{prop:constraint}. Item (3) is Proposition \ref{prop:uniqueness}. Item (4) is Proposition \ref{prop:factor} together with Proposition \ref{prop:benchmark}. Item (5) is Proposition \ref{prop:minimal}. These items show both that the collapse is strict and that the smaller nucleus is still the full benchmark-facing datum at current bridge-object scope.
\end{proof}

\begin{corollary}[What remains open after this theorem]
\label{cor:next}
After Theorem \ref{thm:main}, the next honest primary manuscript below the current bridge-object frontier must do at least one of the following:
\begin{enumerate}
    \item derive or justify the minimal exact-shell affine nucleus from explicit substrate structure in a way that changes benchmark-facing status rather than merely relocating the unexplained layer deeper;
    \item identify a genuinely smaller theorem-level core strictly inside \(\AffineNucleus{\sigma}\) and prove a still sharper collapse to it;
    \item do either of the previous two while preserving the same benchmark one-metric package, the same ordinary matter route, and the same claim boundary.
\end{enumerate}
It is no longer enough merely to restate the former larger primitive reversible constrained dynamics normal form, to insert another deeper same-output relay below that larger package, or to reopen older screened layers as if they were still the live burden.
\end{corollary}

\begin{proof}
Theorem \ref{thm:main} spends the honest internal-compression option that remained open inside the previously live bridge object. Once the larger full normal form has collapsed to the smaller exact-shell affine nucleus, any later manuscript that merely re-expands the larger constraint presentation or merely inserts another same-output deeper relay below it no longer changes benchmark-facing status. What remains is to force the smaller nucleus itself, or compress it further.
\end{proof}

\section{Conclusion}

The positive result of this note is narrower than a completed first-principles substrate derivation and stronger than another deeper same-output relay. The previous bridge-object forcing theorem had already shown that the primitive reversible constrained substrate dynamics normal form is forced by an ambient shell law and is unique there up to the translated-shell and orthogonal-hidden coordinate identifications. The present note then sharpens that result by showing that the larger full normal form still contains redundant constraint-presentation data. At benchmark scope, the live bridge object collapses further to a smaller minimal exact-shell affine nucleus.

That gain is exact and disciplined. Once the shell-coordinate body, shell-data readout, exact shell affine plane, exact shell realization, off-shell gap package, symmetry / time-reversal structure, and solvability / coercive package are fixed, the explicit constraint-target space, linear constraint map, and fixed constraint value are no longer independently benchmark-facing freedom. They are induced canonically by orthogonal projection to the shell-normal complement. All current benchmark-facing downstream uses of the larger full normal form already factor through this smaller nucleus, including the explicit substrate package, the unique selector, and the same one-metric observer package. The benchmark boundary therefore remains unchanged throughout: the one operative metric is still exactly \(\CompletedMetric\), the ordinary matter route is unchanged, there is no second operative metric, and the low-energy matter law is not enlarged.

The open burden is consequently cleaner. The next benchmark-facing manuscript should derive or justify the exact-shell affine nucleus from explicit substrate structure, unless the sharper honest gain available instead is a still smaller theorem-level collapse strictly inside that nucleus. Until then, adoption and derivation should continue to be kept sharply separated. The present note should be read for what it is: a genuine internal bridge-object compression theorem on the active line of \TheoryName{}, not a claim that the full first-principles program is finished.

\end{document}
